% Clase 2 --- Geometría de la integral de Coulomb para una distribución
% continua (§5.3). Descifra la notación r / r' / (r - r') de la fórmula:
% r  = posición del punto de observación (donde está q0),
% r' = posición del elemento de carga dq = rho(r') dV',
% r - r' = vector que va del elemento a q0 (dirección de dF).
\begin{center}
\begin{tikzpicture}[scale=1]

  % ---- volumen cargado ----
  \draw[figline, fill=figblue!7]
    plot[smooth cycle, tension=0.75]
    coordinates {(-0.4,1.25) (1.15,1.55) (2.1,0.7) (1.75,-0.55) (0.4,-1.0) (-0.75,-0.15)};
  \node[figlbl, figblue] at (0.15,1.95) {$V$, \ $\rho(\vec r\,')$};

  % ---- elemento de volumen ----
  \fill[figred] (1.05,0.35) rectangle (1.35,0.65);
  \node[figlbl, figred, above right=0pt] at (1.35,0.62) {$dq = \rho(\vec r\,')\,dV'$};

  % ---- origen ----
  \node[figneutra, minimum size=4pt] (O) at (-2.3,-1.6) {};
  \node[figlbl, below left=0pt] at (O) {$O$};
  \draw[figaxis] (-2.3,-1.6) -- (-1.3,-1.6);
  \draw[figaxis] (-2.3,-1.6) -- (-2.3,-0.6);

  % ---- carga de prueba ----
  \node[figpos] (q0) at (4.4,-1.1) {};
  \node[figlbl, above=4pt] at (q0) {$q_0$};

  % ---- vectores r' y r ----
  \draw[figvecaux] (O) -- (1.2,0.5);
  \node[figlblaux, above left=0pt] at (-0.6,-0.6) {$\vec r\,'$};
  \draw[figvecaux] (O) -- (q0);
  \node[figlblaux, below=2pt] at (1.1,-1.55) {$\vec r$};

  % ---- vector r - r' ----
  \draw[figvec] (1.2,0.5) -- (4.2,-1.05);
  \node[figlbl, figblue, above right=1pt] at (2.6,-0.2) {$\vec r - \vec r\,'$};

  % ---- contribución dF ----
  \draw[figvecres] (q0) -- ++(-27.3:1.3);
  \node[figlbl, figred, below=2pt] at (5.2,-1.5) {$d\vec F$};

\end{tikzpicture}\\[0.5em]
{\small\itshape Cada elemento $dV'$ se trata como una carga puntual y aporta
$d\vec F$ según Coulomb; la suma pasa a ser una integral sobre $V$. La fórmula
sólo dice eso: $|\vec r - \vec r\,'|$ es la distancia entre el elemento y el
punto de observación, y $(\vec r - \vec r\,')/|\vec r - \vec r\,'|$ es el versor
que va \textbf{del elemento hacia $q_0$}. No se supuso $\rho$ uniforme: puede
variar punto a punto e incluso cambiar de signo.}
\end{center}
